% !TeX root = ../main.tex
\chapter{Testing}

Testing for Jan-Sunwai AI was designed as a layered strategy combining automated checks, manual UAT scripts, resilience drills, and security/load validation. Because the platform supports governance-sensitive workflows, testing emphasized not only feature correctness but also operational integrity under degraded conditions [18, 7, 8].

\section{Testing Plan}

The testing plan was structured to validate correctness, security, and operational resilience across the full grievance lifecycle. Instead of relying on a single test phase, verification activities were layered and repeated at key implementation checkpoints to reduce late-stage defects.

\subsection{Testing Objectives}

Testing objectives were scoped to cover the distinct failure modes that a production civic system must defend against:

\begin{itemize}
\item Validate correctness of core APIs and role authorization boundaries.
\item Verify AI-path behavior for valid, invalid, and failure scenarios.
\item Confirm notification chain consistency across complaint lifecycle changes.
\item Validate resilience behavior under model/database disruptions.
\item Ensure production-readiness through integration, load, and UAT checks.
\end{itemize}

\subsection{Test Environment}

The testing environment baseline outlined below ensures consistency across all validation phases.


\begin{description}
\item[Backend runtime:] FastAPI application with test clients and dependency overrides.
\item[Database mode:] MongoDB-backed flows with selective mock/fake collection tests.
\item[Frontend mode:] Vite build/lint validation and persona-driven manual walkthroughs.
\item[Model runtime mode:] Ollama available for normal flows; manually disabled for resilience checks.
\item[Automation tools:] \texttt{pytest}, \texttt{httpx}, \texttt{locust}, and script wrappers for platform-specific execution.
\end{description}

\subsection{Coverage Matrix}

The coverage matrix below maps each test suite to its type and the specific system behaviours it validates:


\begin{xltabular}{\textwidth}{|>{\RaggedRight\arraybackslash}p{0.30\textwidth}|>{\RaggedRight\arraybackslash}p{0.24\textwidth}|>{\RaggedRight\arraybackslash}X|}
\hline
\textbf{Test Artifact} & \textbf{Type} & \textbf{Coverage Highlights} \\ \hline
\endhead
\path{test_schemas.py} & Unit Validation & GeoLocation bounds, complaint schema validity, confidence-score guardrails \\ \hline
\path{test_auth_permissions.py} & Unit Authorization & Token claims, admin and worker/dept-head access guards \\ \hline
\path{test_api_integration.py} & Integration Smoke & Root route, health checks, versioned aliases, lightweight CI path \\ \hline
\path{test_api_matrix.py} & Integration Matrix & /api/v1 profile route, CORS preflight behavior, security header presence \\ \hline
\path{test_notification_chain.py} & Integration Behavior & Status update notification creation, email hook path, mark-all-read logic \\ \hline
\path{test_resilience_security.py} & Resilience + Security & Upload hardening, analyze degrade 503, readiness degradation, sanitization \\ \hline
Manual UAT scripts & User Acceptance & Citizen and admin journeys, friction logging, UX verification \\ \hline
Load scripts (locust) & Performance & Endpoint mix throughput and graceful degradation behavior \\ \hline
\caption{Testing Coverage Matrix by Suite}
\end{xltabular}

\section{Types of Testing}

Multiple testing types were applied because the platform combines API logic, role-sensitive workflows, AI-assisted features, and deployment constraints. This blended strategy ensured both component-level reliability and end-to-end behavioral confidence.

\subsection{Unit Testing}

Unit tests focused on isolated validation and control logic.

\begin{table}[H]
\centering
\begin{tabularx}{\textwidth}{|>{\RaggedRight\arraybackslash}X|>{\RaggedRight\arraybackslash}X|>{\RaggedRight\arraybackslash}X|}
\hline
\textbf{Case ID} & \textbf{Area} & \textbf{Expected Behavior} \\ \hline
UT-01 & GeoLocation bounds & Reject latitude outside $[-90, 90]$ and longitude outside $[-180, 180]$ \\ \hline
UT-02 & AI confidence schema & Reject confidence values outside $[0, 1]$ \\ \hline
UT-03 & Admin role check & Non-admin caller receives HTTP 403 on admin-only dependency \\ \hline
UT-04 & Worker approval gate & Unapproved worker cannot pass worker-only dependency path \\ \hline
UT-05 & Sanitization function & Script payload transformed to escaped safe representation \\ \hline
\end{tabularx}
\caption{Representative Unit Test Cases}
\end{table}

\subsection{Integration Testing}

Integration tests verified route composition, middleware behavior, and cross-router guarantees such as versioned aliases and headers.

\begin{itemize}
\item Root and liveness endpoints confirmed service online state.
\item Versioned alias endpoints under \texttt{/api/v1} validated integration compatibility.
\item Security headers confirmed on endpoint responses.
\item Preflight CORS behavior validated for allowed origins.
\end{itemize}

\subsection{System Testing}

System testing executed end-to-end flows across authentication, analyze, complaint lifecycle, assignment, and notification propagation. This stage focused on behavior continuity rather than isolated endpoint correctness.

\begin{itemize}
\item \textbf{ST-01 (Citizen flow):} Upload $\rightarrow$ analyze $\rightarrow$ submit. Ensures the complaint record is created with its status history initialized.
\item \textbf{ST-02 (Dept-head update):} Dept-head updates status. Ensures the citizen receives a notification and the timeline reflects the transition.
\item \textbf{ST-03 (Worker completion):} Worker marks a complaint as done. Ensures the complaint transitions and worker slot logic updates as expected.
\item \textbf{ST-04 (Admin triage):} Admin triage decision. Ensures an ambiguous complaint successfully completes its governed label decision path.
\item \textbf{ST-05 (Data export):} CSV export generation. Ensures the export file is generated with expected schema columns.
\end{itemize}

\subsection{User Acceptance Testing (UAT)}

The UAT plan was executed through persona scripts:

\begin{itemize}
\item Citizen script: login, analyze, map review, submit, dashboard verification.
\item Admin script: queue filtering, bulk actions, heatmap view, reassignment, CSV export.
\item Feedback loop: friction capture with severity ranking and follow-up fixes.
\item Resilience check: AI service interruption and recovery behavior in UI.
\end{itemize}

\subsection{Performance and Load Testing}

Load testing on baseline server configuration (48GB RAM, H100 GPU) over 30 minutes with 50 concurrent civic workers confirmed responsive API performance and graceful degradation at peak load. The deterministic routing threshold successfully captures ambiguous images with high accuracy for auto-routing.

\section{Testing Techniques}

Testing combined behavioral, negative, resilience, and load techniques to ensure stability. \textbf{Black-Box Testing} validated API contracts from a consumer perspective. \textbf{Negative Testing} deliberately exercised invalid tokens, corrupted uploads, and unauthorized role access. \textbf{Resilience Testing} simulated database unreachability and Ollama service failures to verify controlled degradation. Finally, \textbf{Load Testing} subjected read-heavy routes and background queues to simulated traffic, confirming graceful fallback mechanisms during peak pressure [18].

\subsection{Defect Logging}

Defects were logged using a severity-and-impact model with explicit reproduction details, observed behavior, expected behavior, and fix status. Regressions were re-run in the relevant unit/integration suite before closure to avoid issue recurrence.

\subsection{System Evaluation}

The platform was validated to support its intended civic workflow, with core CRUD application remaining highly responsive during background AI processing. Fallback and escalation paths ensure service availability under model unavailability.


\section{Testing Outcomes}

\begin{table}[H]
\centering
\begin{tabularx}{\textwidth}{|>{\RaggedRight\arraybackslash}X|>{\RaggedRight\arraybackslash}X|>{\RaggedRight\arraybackslash}X|}
\hline
\textbf{Validation Area} & \textbf{Outcome} & \textbf{Notes} \\ \hline
Backend automated suite & Pass & Core suite executed successfully with latest observed run showing full pass status \\ \hline
Security checks & Pass & Upload hardening, sanitization, and authorization boundaries validated \\ \hline
Resilience checks & Pass & Controlled degraded responses confirmed for model/DB disruption paths \\ \hline
Frontend lint/build & Pass & UI changes and route additions validated through build and lint cycle \\ \hline
Compose validation & Pass & Unified compose configuration validated for default and prod profile flows \\ \hline
UAT persona scripts & Pass with refinements & Minor UX adjustments applied in same sprint feedback loop \\ \hline
\end{tabularx}
\caption{Outcome Summary by Validation Area}
\end{table}

\section{Residual Risks and Test Gaps}

Although major paths were validated, future strengthening can include:

\begin{itemize}
\item Long-duration soak tests on production-equivalent infrastructure.
\item Expanded browser matrix for UI behavior consistency.
\item External vulnerability scanning integration in release pipeline.
\item Higher-volume synthetic image corpus evaluation for classifier drift tracking.
\end{itemize}

\section{Chapter Summary}

Testing combined automated rigor with practical operational validation. The resulting quality posture supports both technical confidence and governance readiness, while documented residual risks provide a clear roadmap for post-submission hardening.

